\subsection{CP phase in modular flavor models and discrete Froggatt-Nielsen models --- arXiv:2312.11809}

\textbf{Authors:} Shota Kikuchi, Tatsuo Kobayashi, Kaito Nasu

\paragraph{主な主張}
固定点近傍の残留電荷で階層を作る模型では、一つの小さな複素パラメータだけで十分なCP破れを得ることが難しく、検討した構成で複数モジュライが両者の実現に有効であると示した\cite{Kikuchi:2023fpl}。

\paragraph{新規性}
モジュラーフレーバー模型と離散Froggatt--Nielsen模型に共通する抑制因子の構造を用い、質量階層とCP位相を両立する多モジュライ質量行列を提示した。

\paragraph{位置付け}
残留対称性から自然な階層を生成する際に、CP破れが追加の条件となることを明らかにする研究である。
